\section{Introduction}
\label{sec:intro}

The pre-train and fine-tune paradigm has fundamentally reshaped modern machine learning, enabling a single backbone to be adapted to a multitude of downstream tasks via parameter-efficient fine-tuning (PEFT) methods like LoRA \cite{hu2021lora, peft2022, lialin2023peft}. As the number of specialized downstream experts grows, ensembling them without retraining has emerged as a key area of interest. Model merging techniques, including task arithmetic \cite{ilharco2022editing, ortiz2020task}, TIES-Merging \cite{yadav2023ties}, RegMean \cite{jin2022regmean}, Fisher Merging \cite{matena2021merging}, and Model Soups \cite{wortsman2022model}, allow practitioners to combine multiple specialized parameter sets into a single multi-task model.

More recently, research has moved toward \emph{dynamic test-time model merging}, where expert weights are merged on-the-fly based on the properties of the incoming query stream. A notable milestone is Parameter-Free Subspace Routing (PFSR) \cite{pfsr}, which computes sample-specific routing coefficients by projecting penultimate representations onto frozen task classifier heads. This approach operates without any trainable router parameters or calibration data, offering an elegant, lightweight solution for test-time dynamic adaptation.

However, a severe and fundamental limitation arises when these parameter-space dynamic merging techniques are deployed in realistic, heterogeneous streaming environments. When an incoming batch contains a mixture of different tasks (e.g., MNIST and CIFAR-10), parameter-space merging methods are forced to average the routing coefficients over the batch dimension to maintain a single set of merged weights for the forward pass. This averaging results in \emph{heterogeneity collapse}: the merged parameters collapse to a static, sub-optimal uniform merge, destroying the individual specialization of each expert and causing accuracy to plunge from high-performing homogeneous baselines to near-uniform levels.

To combat heterogeneity collapse, state-of-the-art methods like Micro-Batch Homogenization (MBH) \cite{mbh} wrap the model in a stateful systems scheduling pipeline. MBH dynamically buffers incoming samples, sorts them by task similarity, and partitions them into homogeneous micro-batches before ensembling. While effective, this systems-centric approach introduces significant bloat:
\begin{itemize}
    \item \textbf{Serving Complexity:} It introduces a stateful scheduling layer in the inference server, breaking the simple stateless paradigm of deep learning inference.
    \item \textbf{Queuing Latency:} It forces query streams to be delayed in buffers, introducing unpredictable queuing latencies that are unacceptable for real-time applications.
    \item \textbf{Memory \& Compute Overhead:} The dynamic grouping, sorting, and partitioning algorithms consume CPU and memory, adding system-level bottleneck dependencies.
\end{itemize}

Guided by the principle of Occam's razor, we ask a fundamental question: \emph{Can we achieve perfect robustness to heterogeneous streaming entirely at the network level, without any stateful buffering or systems scheduling wrappers?}

We answer in the affirmative. In this work, we present \textbf{SABLE (Sample-wise Activation Blending of Low-Rank Experts)}, a minimalist and mathematically elegant solution that completely bypasses the need for weight-space averaging. Our core insight is that ensembling does not have to happen in weight space. By leveraging the distributive property of matrix multiplication, we shift the ensembling step from parameter space to activation space. SABLE executes a shared forward pass through the pre-trained base model, alongside parallel, extremely lightweight low-rank ($r=8$) adapter passes, blending their activations on-the-fly using sample-wise coefficients derived from non-parametric cosine subspace projections.

SABLE possesses several compelling advantages:
\begin{enumerate}
    \item \textbf{Perfect Heterogeneity Robustness:} By ensembling activations on a per-sample basis, SABLE is natively immune to heterogeneity collapse. It maintains identical, robust generalization performance under both homogeneous and fully mixed task streams.
    \item \textbf{System-Agnostic Simplicity:} SABLE strips away the entire stateful scheduling, buffering, sorting, and partitioning pipeline of MBH. It can be deployed in standard, stateless serving frameworks with zero modifications to the systems stack.
    \item \textbf{Extreme Resource Efficiency:} By parameterizing specialized experts as low-rank LoRA adapters ($r \ll D$), SABLE minimizes parameter storage. The activation-space blending layer is extremely lightweight, introducing negligible compute overhead.
\end{enumerate}

We evaluate SABLE in a synthetic 14-layer, 192-dimensional Analytical Coordinate Sandbox simulating multi-task streams. SABLE's default Late Adaptation configuration achieves a flatline joint mean accuracy of \textbf{68.10\%} across both homogeneous and heterogeneous streams, completely avoiding heterogeneity collapse (0.00\% collapse) and outperforming the state-of-the-art, complex systems-level pipeline PFSR+MBH (\textbf{67.20\%}) natively. SABLE achieves this superior performance while entirely eliminating systems-level latency, queue buffering, or state dependencies, returning model serving to its clean, stateless, and highly reproducible roots.
